\newcommand{\NoLink}[1]{{\protect\NoHyper#1\protect\endNoHyper}} % Allows the use of \citeauthor without actually referencing back to the reference list, in case you want to just say the author's name. Example: \NoLink{\citeauthor{CoolKid2019}}
\newcommand\invisiblechapter[1]{                % makes an invisible chapter!
  \refstepcounter{chapter}%
  \addcontentsline{toc}{chapter}{\protect\numberline{\thechapter}#1}%
  \chaptermark{#1}}
\newcommand\invisiblesection[1]{                % makes an invisible section!
  \refstepcounter{section}%
  \addcontentsline{toc}{section}{\protect\numberline{\thesection}#1}%
  \sectionmark{#1}}
\newcommand\invisiblesubsection[1]{%            % makes an invisible subsection!
  \refstepcounter{subsection}%
  \addcontentsline{toc}{subsection}{\protect\numberline{\thesubsection}#1}%
  \subsectionmark{#1}}

\newcommand\diff[2][]{\mathrm{d}^{#1}{#2}}      % creates a nth-order ordinary differential variable
\newcommand\pdiff[2][]{\partial^{#1}{#2}}       % creates a nth-order partial differential variable
\newcommand\sdiff[2][]{\mathrm{D}^{#1}{#2}}     % creates a nth-order substantial differential variable
\newcommand\deriv[3][]{                         % creates a nth-order ordinary derivative
\frac{\diff[#1]{#2}}{\diff{#3}^{#1}}
}
\newcommand\pderiv[3][]{                        % creates a nth-order partial derivative
\frac{\pdiff[#1]{#2}}{\pdiff{#3}^{#1}}
}
\newcommand\sderiv[3][]{                        % creates a nth-order substantial derivative
\frac{\sdiff[#1]{#2}}{\sdiff{#3}^{#1}}
}
\newcommand\cderiv[3][]{\diff[#1]{}_{#3}{#2}}   % creates a compact form nth-order ordinary derivative
\newcommand\cpderiv[3][]{\pdiff[#1]{}_{#3}{#2}} % creates a compact form nth-order partial derivative
\newcommand\csderiv[3][]{\sdiff[#1]{}_{#3}{#2}} % creates a compact form nth-order substantial derivative

\newcommand\uVec[2]{\hat{\Vec{#1}}_{#2}}        % creates a unit vector
\newcommand\uVecT[2]{\overline{\Vec{#1}}_{#2}}  % creates a transformed unit vector
\newcommand\transpose{^\textsf{T}}              % places in a transpose symbol

%https://tex.stackexchange.com/questions/417310/double-bar-notation-for-second-order-tensors4
\newcommand\barbelow[1]{\stackunder[1.6pt]{$#1$}% creates a bar below an object
                        {\rule{0.9ex}{.1ex}}}   

\newcommand{\grad}{ \barbelow{\nabla}}          % \grad now includes bar below.
\newcommand{\curl}{ \barbelow{\nabla}\times}    % \curl now includes bar below.
\newcommand\proj[2]{ \mathsf{proj}_{#1}^{#2}\,} % Projection operator
\newcommand{\Res}{ \mathsf{Res}}                % Residual operator

%% Create typical dimensionless numbers used in fluid + heat transfer problems
\renewcommand{\Pr}{\mathrm{Pr}}                 % Prandtl number
\renewcommand{\Re}{\mathrm{Re}}                 % Reynolds number
\newcommand{\Fr}{\mathrm{Fr}}                   % Froude number
\newcommand{\Eu}{\mathrm{Eu}}                   % Euler number
\newcommand{\St}{\mathrm{St}}                   % Strouhal number
\newcommand{\Gr}{\mathrm{Gr}}                   % Grashof number
\newcommand{\Ma}{\mathrm{Ma}}                   % Mach number
\newcommand{\Pe}{\mathrm{Pe}}                   % Péclet number
\newcommand{\Tu}{\mathrm{Tu}}                   % Turbulence level

%% With Eulervm, no \therefore defined. newtxmath provides \therefore, so this is only defined if missing.
% https://tex.stackexchange.com/questions/142551/how-to-put-therefore-and-implies-symbols
\providecommand{\therefore}{\dot{.\hspace{.055in}.}\hspace{.2in}}

%% standard /mapsto symbol, but reversed
% https://tex.stackexchange.com/questions/26508/left-version-of-mapsto
\providecommand{\mapsfrom}{\mathrel{\reflectbox{\ensuremath{\mapsto}}}}
\providecommand{\longmapsfrom}{\mathrel{\reflectbox{\ensuremath{\longmapsto}}}}

%% Using the same figure twice with no new number
% https://tex.stackexchange.com/questions/200211/using-the-same-figure-twice-with-no-new-number
\newcommand{\repeatcaption}[2]{%
  \renewcommand{\thefigure}{\ref{#1}}%
  \captionsetup{list=no}%
  \caption{#2 (repeated from page \pageref{#1})}%
  \addtocounter{figure}{-1}% So that next figure after the repeat gets the right number.
}

% Create delimiter environments (e.g floor(t) produces ⌊t⌋)
% --------------------
% https://tex.stackexchange.com/questions/118173/how-to-write-ceil-and-floor-in-latex
\DeclarePairedDelimiter\ceil{\lceil}{\rceil}    % Use ceil*{...} to produce big brackets
\DeclarePairedDelimiter\floor{\lfloor}{\rfloor} % Use floor*{...} to produce big brackets

% Create new math operators
% --------------------
\DeclareMathOperator*{\argmax}{arg\,max}
\DeclareMathOperator*{\argmin}{arg\,min}

% tensor - double dot product
\newcommand{\cddot}{\mathbin{{\cdot}{\cdot}}}

% Renew Commands:
% --------------------
\renewcommand{\vec}[1]{ \barbelow{\bm{#1}} }    % \vec to \bm combination
\renewcommand{\div}{ \barbelow{\nabla}\cdot}    % \div now includes bar below.

\renewcommand*{\tocheadstart}{\vspace*{0pt}}    % straight from the memoir manual, may need to change if you change
\renewcommand*{\lofheadstart}{\vspace*{0pt}}    % titlespacing for chapters... investigate?
\renewcommand*{\lotheadstart}{\vspace*{0pt}}    % pre-title spacing is now the same for toc, lof, lot and nomenclature

% \itemize labels %% https://texblog.org/2008/10/16/lists-enumerate-itemize-description-and-how-to-change-them/
\renewcommand{\labelitemi}{-}
\renewcommand{\labelitemii}{-}
\renewcommand{\labelitemiii}{-}
\renewcommand{\labelitemiv}{-}

% Robust Commands:
% --------------------
\DeclareRobustCommand{\VAN}[0]{Van} % Within text, authors with a van/Van prefix use Van within text, van outside text. 
